\chapter{引言}

在网络化的多智能体系统中，分布式协作优化是一个不可或缺的核心任务。在理想的无对抗网络环境中，多智能体优化的标准范式是假设网络中的每个智能体持有一个局部代价函数，随后所有的智能体通过点对点的局部通信相互协作，共同寻找一个能够最小化全体局部代价函数算术平均值的全局最优解\upcite{nedic2009distributed}。这一基于算术平均的协同寻优框架在过去十几年中推动了传感器网络、分布式控制等领域的蓬勃发展，并在理论上得到了充分且成熟的研究。

然而，现实的分布式计算环境很容易遭受恶意攻击，现有的基于算术平均值目标的多智能体优化范式在面临拜占庭故障时十分脆弱。在拜占庭模型中\upcite{lynch1996distributed}，未知子集的智能体可能表现出任意的对抗性行为，甚至向不同的邻居发送自相矛盾的虚假信息。由于算术平均值对极端异常数据几乎没有任何抵抗力，因此即使网络中仅存在单一的拜占庭节点，它也能完全劫持并摧毁整个多智能体网络的全局优化目标\upcite{sundaram2019distributed}。此外，现有的容错一致性算法也无法处理梯度寻优问题\upcite{vaidya2012iterative}，而容错机器学习算法又高度依赖数据独立同分布假设以及中心化服务器\upcite{yin2018byzantine}，导致前人方法均无法直接解决无中心、数据非独立同分布场景下的拜占庭优化难题。

针对上述问题，本文在无中心协调器的前提下，提出了一种全新的拜占庭容错多智能体优化问题框架。为了减少恶意节点对优化结果的干扰，我们摒弃了传统的算术平均值目标，在数学上将安全的全局优化目标重新界定为仅由正常节点局部代价函数构成的凸组合。更为关键的是，为了防止该凸组合坍缩为只受极少数正常节点主导的结果，我们创新性地引入了($\beta, \gamma$)-可容许性这一严格的数学度量，明确要求最终的全局决策必须至少包含$\gamma$个正常节点的信息，且被采纳的每个正常节点数据的权重占比都不得低于$\beta$。

在确立了该合法目标的结构性约束后，我们进一步从理论底层探索了该问题的数学极限，并设计了相应的求解算法。首先，通过不可区分性论证，我们证明了在允许拜占庭节点发送不一致消息的网络中，存在一个不可逾越的不可能性边界。其次，为了逼近这一边界，我们提出了一种基于双重修剪的分布式梯度下降算法。该算法摒弃了单一的状态聚合，转而在局部状态估计和局部梯度计算两个维度上同时引入基于排序的修剪过滤机制，从而有效抵御了多维度的对抗性攻击。

理论分析与收敛性证明表明，本文提出的算法能够确保所有正常智能体的局部估计值渐近收敛到满足($\beta, \gamma$)-可容许性的合法最优解$\mathcal{X}(\beta, \gamma)$集合中。特别地，当底层通信网络为完全图时，该算法所能保证的$\gamma$值完美匹配了不可能性定理中的理论上限，且其实现的$\beta$值达到了理论最优值的$\frac{1}{2}$。这标志着在最恶劣的拜占庭环境下，系统依然能在理论极限附近守住正常节点的协同权益。

综上所述，本文的核心贡献在于：
\begin{itemize}
    \item 提出了一种无中心拜占庭容错多智能体优化框架，并首创了($\beta, \gamma$)-可容许性度量，以严格刻画和约束容错全局目标的凸系数结构。
    \item 通过证明该优化问题的不可能性定理，确立了任何分布式算法在防范拜占庭攻击时所能达到的理论物理上限。
    \item 提出了一种基于双重修剪的分布式梯度下降算法，并严格证明了其在完全图网络中能够达到与理论极限相匹配的最优容错性能。
\end{itemize}
